\PassOptionsToPackage{unicode}{hyperref}
\PassOptionsToPackage{hyphens}{url}
\documentclass[12pt,a4paper]{report}
\usepackage[a4paper,margin=1in]{geometry}
\usepackage{amsmath,amssymb}
\usepackage{iftex}
\ifPDFTeX
  \usepackage[T1]{fontenc}
  \usepackage[utf8]{inputenc}
  \usepackage{textcomp} % provide euro and other symbols
\else % if luatex or xetex
  \usepackage{unicode-math} % this also loads fontspec
  \defaultfontfeatures{Scale=MatchLowercase}
  \defaultfontfeatures[\rmfamily]{Ligatures=TeX,Scale=1}
\fi
\usepackage{lmodern}
\ifPDFTeX\else
\fi
\IfFileExists{upquote.sty}{\usepackage{upquote}}{}
\IfFileExists{microtype.sty}{% use microtype if available
  \usepackage[]{microtype}
  \UseMicrotypeSet[protrusion]{basicmath} % disable protrusion for tt fonts
}{}
\makeatletter
\@ifundefined{KOMAClassName}{% if non-KOMA class
  \IfFileExists{parskip.sty}{%
    \usepackage{parskip}
  }{% else
    \setlength{\parindent}{0pt}
    \setlength{\parskip}{6pt plus 2pt minus 1pt}}
}{% if KOMA class
  \KOMAoptions{parskip=half}}
\makeatother
\usepackage{xcolor}
\usepackage{graphicx}
\usepackage{longtable,booktabs,array}
\usepackage{calc} % for calculating minipage widths
\usepackage{etoolbox}
\makeatletter
\patchcmd\longtable{\par}{\if@noskipsec\mbox{}\fi\par}{}{}
\makeatother
\IfFileExists{footnotehyper.sty}{\usepackage{footnotehyper}}{\usepackage{footnote}}
\makesavenoteenv{longtable}
\setlength{\emergencystretch}{3em} % prevent overfull lines
\providecommand{\tightlist}{%
  \setlength{\itemsep}{0pt}\setlength{\parskip}{0pt}}
\setcounter{secnumdepth}{-\maxdimen} % remove section numbering
\ifLuaTeX
  \usepackage{selnolig}  % disable illegal ligatures
\fi
\usepackage{bookmark}
\IfFileExists{xurl.sty}{\usepackage{xurl}}{} % add URL line breaks if available
\urlstyle{same}
\hypersetup{
  hidelinks,
  pdfcreator={LaTeX via pandoc}}

\author{}
\date{}
\usepackage{tocloft}
\setlength{\cftbeforesecskip}{0pt}
\setlength{\cftbeforesubsecskip}{0pt}
\setlength{\cftbeforesubsubsecskip}{0pt}
\setlength{\cftbeforechapskip}{0pt}

% --- Report conversion: print "Chapter N" above each main heading ---
% (section numbering stays disabled, so headings keep their manual numbers
%  such as "1.1 Document Purpose"; the label is supplied via \chapterlabel)
\makeatletter
\newcommand{\chapterlabel}{}
\def\@makechapterhead#1{%
  \vspace*{50\p@}%
  {\parindent \z@ \raggedright \normalfont
    \ifx\chapterlabel\@empty\else
      \huge\bfseries \chapterlabel \par\nobreak
      \vskip 20\p@
    \fi
    \interlinepenalty\@M
    \Huge \bfseries #1\par\nobreak
    \vskip 40\p@
  }}
\makeatother

\begin{document}
\begin{titlepage}
 \centering
{\scshape\bfseries\LARGE SOFTWARE REQUIREMENT SPECIFICATION
(SRS) \par}
\vspace{1cm}
{\scshape\bfseries\huge COMPUTER VISION AND AI-BASED CONSTRUCTION SITE MONITORING VIA SATELLITE IMAGERY\par}
\vspace{1cm}
{\scshape\bfseries\huge Version 1.0 \par}
\vspace{1.2cm}
{\huge\bfseries Usman Rasheed Siddiqui (CS-24038) \par}
\vspace{1.2cm}
{\huge\bfseries Salman Raza (CS-24043) \par}
\vspace{1.2cm}
{\huge\bfseries Ahsan Atif Rajput (CS-24044) \par}
\vspace{1.2cm}
{\huge\bfseries Danish Farooqui (CS-24070) \par}
\vspace{1cm}
{\Large\itshape Department of Computer and Information Systems
\\ NED University of Engineering and Technology\par}
\vfill
Submitted to\par
Ms. Fakhra Aftab \textsc
\vfill
% Bottom of the page
{\large October 20, 2025 \par}
\end{titlepage}
\newpage

{
\setcounter{tocdepth}{3}
\tableofcontents
}
\newpage

\vspace{5cm}
\renewcommand{\chapterlabel}{Chapter 1}
\chapter{1. Introduction}\label{introduction}

\section{1.1 Document Purpose}\label{document-purpose}

This Software Requirements Specification (SRS) describes the functional
and non-functional requirements for the Computer Vision and AI-Based
Construction Site Monitoring via Satellite Imagery system. It defines
the expected system behaviour, interfaces, data requirements,
constraints, and interaction models to guide design, implementation, and
verification.

\section{1.2 Product Scope}\label{product-scope}

The \textbf{Computer Vision and AI-Based Construction Site Monitoring
via Satellite Imagery} system aims to monitor construction progress by
detecting physical changes in construction sites using satellite
imagery, computer vision, and artificial intelligence. The system will
allow authorized users, including builders, owners, and site
supervisors, to register construction sites by providing the site name,
geographical coordinates, and construction start date. Using Google
Earth Engine, the system will retrieve available Sentinel-2 satellite
images for the selected location and date range, preprocess them, and
compare images captured at different times to identify and highlight
changed areas.

The system will display detected changes, estimated changed areas, and
AI-generated insights through a user-friendly dashboard. It will display
the current detected changes for a selected monitoring check and allow
users to generate and download reports. The project focuses on
\textbf{site-level change detection and progress monitoring} and will
not determine exact construction completion percentages, assess
structural safety, or replace physical site inspections. Results will
depend on image resolution, availability, cloud cover, and model
accuracy.

\section{1.3 Intended Audience and Document
Overview}\label{intended-audience-and-document-overview}

This document is intended for the project team, software and AI
developers, reviewers, instructors, and other stakeholders who need to
understand the system requirements. Chapter 1 introduces the project;
Chapter 2 describes the product context and stakeholders; Chapter 3
specifies interfaces, user and system functional requirements, and
behavioural models; Chapter 4 defines non-functional requirements;
Chapter 5 records any additional requirements; and the appendices
describe the data dictionary and traceability information.

\section{1.4 Definitions, Acronyms and
Abbreviations}\label{definitions-acronyms-and-abbreviations}

\begin{longtable}[]{@{}
  >{\raggedright\arraybackslash}p{(\columnwidth - 2\tabcolsep) * \real{0.5000}}
  >{\raggedright\arraybackslash}p{(\columnwidth - 2\tabcolsep) * \real{0.5000}}@{}}
\toprule\noalign{}
\begin{minipage}[b]{\linewidth}\raggedright
Term
\end{minipage} & \begin{minipage}[b]{\linewidth}\raggedright
Definition
\end{minipage} \\
\midrule\noalign{}
\endhead
\bottomrule\noalign{}
\endlastfoot
AI & Artificial Intelligence \\
CV & Computer Vision \\
DFD & Data Flow Diagram \\
FR & Functional Requirement \\
GEE & Google Earth Engine \\
IoU & Intersection over Union \\
NFR & Non-Functional Requirement \\
SFR & System Functional Requirement \\
SRS & Software Requirements Specification \\
Sentinel-2 & Earth-observation satellite mission providing multispectral
imagery \\
UFR & User Functional Requirement \\
\end{longtable}

\section{1.5 Document Conventions}\label{document-conventions}

Requirements are identified by stable IDs where available. ``Shall''
indicates a mandatory requirement. Tables are used to describe detailed
system requirements, including function, inputs, outputs, dependencies,
actions, preconditions, postconditions, and side effects. Requirement
IDs from the supplied report are preserved in this draft;
inconsistencies and gaps should be reviewed before final submission.

\section{1.6 References and
Acknowledgments}\label{references-and-acknowledgments}

\begin{enumerate}
\def\labelenumi{\arabic{enumi}.}
\item
  Python 3.11+ Documentation
\item
  Django Documentation
\item
  Supabase Documentation
\item
  Google Earth Engine Documentation
\item
  Copernicus Sentinel-2 Satellite Imagery Documentation
\item
  StarUML --- UML Modeling Tool used to create Use Case Diagrams,
  Sequence Diagrams, and Data Flow Diagrams (DFDs).
\item
  ChatGPT --- AI tool used for assistance in SRS structuring,
  documentation, and requirement refinement.
\item
  Google Gemini --- AI tool used for assistance in research, idea
  development, and documentation.
\item
  Class Diagram and Use Case Diagram --- Reference material.
\item
  Software Requirements Specification (SRS) Template --- Reference
  structure provided for the project.
\end{enumerate}

\renewcommand{\chapterlabel}{Chapter 2}
\chapter{2. Overall Description}\label{overall-description}

\section{2.1 Product Perspective}\label{product-perspective}

The system is an internet-based application for site-level construction
change monitoring. Users register a site using a name, bounding-box
coordinates, and construction start date. The system requests suitable
Sentinel-2 imagery through Google Earth Engine, compares imagery from
different dates, generates a change mask and estimated changed-area
measurements, presents AI-generated observations, and supports
monitoring report generation and download.

The system depends on external satellite imagery and geospatial
services. Results depend on image availability, spatial resolution,
cloud cover, image quality, and model performance. The system does not
determine exact construction completion percentages, assess structural
safety, or replace physical site inspections.

\section{2.2 Product Functionality}\label{product-functionality}

\begin{itemize}
\item
  Account registration, login, and logout.
\item
  Construction site registration, coordinate validation, map
  confirmation, and site details viewing.
\item
  Site access requests and creator approval or rejection.
\item
  Retrieval of suitable Sentinel-2 imagery through Google Earth Engine.
\item
  Image preparation, change detection, generation of a change mask, and
  calculation of estimated changed area.
\item
  Storage and retrieval of monitoring records and result references.
\item
  Dashboard display of sites, monitoring status, imagery, detected
  changes, measurements, and AI observations.
\item
  Generation and download of monitoring reports.
\item
  Notifications for invalid input, unavailable imagery, processing
  failures, and external-service or database errors.
\end{itemize}

\section{2.3 Users and
Characteristics}\label{users-and-characteristics}

Direct users are builders, construction site owners, and site
supervisors. The supplied report states that registered users have the
same general application permissions, while a site creator may approve
or reject access requests for their own site.

Indirect stakeholders include the Software Engineer, AI Engineer, and
System Administrator. Domain stakeholders include Google Earth Engine
and the Sentinel-2 satellite imagery provider.

\section{2.4 Operating Environment}\label{operating-environment}

The application is intended to be accessed through a modern web browser
over the internet. The backend requires access to a database, Google
Earth Engine services, and a processing environment capable of accessing
satellite image pixels and running the selected computer-vision method.
Exact deployment operating system, browser versions, memory, and storage
minimums should be confirmed against the implementation.

\section{2.5 Design and Implementation
Constraints}\label{design-and-implementation-constraints}

\begin{itemize}
\item
  Satellite image availability varies by location, date range, cloud
  cover, and image quality.
\item
  Sentinel-2 imagery has finite spatial resolution; results must not
  claim finer spatial detail than the source data supports.
\item
  Google Earth Engine usage is subject to its access, quota, and service
  policies.
\item
  The processing pipeline must access actual image pixel data; storing
  only image references is insufficient for model inference.
\item
  Generated masks and other output files may be stored separately from
  database metadata, with references recorded in the database.
\item
  The application requires internet connectivity for its web interface
  and external imagery services.
\end{itemize}

\section{2.6 User Documentation}\label{user-documentation}

The final system should provide brief guidance for account creation,
site registration, coordinate entry and confirmation, running an image
comparison, interpreting change measurements and AI observations,
viewing saved results, and generating/downloading reports. The actual
user guide can be attached or linked when prepared.

\section{2.7 Assumptions and
Dependencies}\label{assumptions-and-dependencies}

\begin{itemize}
\item
  Users provide site coordinates and a construction start date; they do
  not upload satellite images.
\item
  Google Earth Engine provides access to suitable Sentinel-2 imagery
  when available.
\item
  The database stores site details, image references and metadata,
  monitoring records, measurements, and result references.
\item
  Image pixels must be available to the computer-vision processing
  pipeline.
\item
  All registered users have the same general application permissions;
  site creators have the site-specific ability to approve or reject
  access requests.
\item
  AI-generated observations may be inaccurate and must remain grounded
  in the available image analysis and measurements.
\end{itemize}

\section{2.8 Stakeholder Identification and
Viewpoints}\label{stakeholder-identification-and-viewpoints}

\section{1. List of Stakeholders}\label{list-of-stakeholders}

\begin{itemize}
\item
  Construction Site Users: Builders, Owners, and Site Supervisors
\item
  Software Engineer
\item
  AI Engineer
\item
  System Administrator
\item
  Google Earth Engine
\item
  Sentinel-2 Satellite Imagery Provider
\end{itemize}

\section{2. Stakeholder Viewpoints}\label{stakeholder-viewpoints}

\begin{longtable}[]{@{}
  >{\raggedright\arraybackslash}p{(\columnwidth - 4\tabcolsep) * \real{0.1804}}
  >{\raggedright\arraybackslash}p{(\columnwidth - 4\tabcolsep) * \real{0.2888}}
  >{\raggedright\arraybackslash}p{(\columnwidth - 4\tabcolsep) * \real{0.5309}}@{}}
\toprule\noalign{}
\endhead
\bottomrule\noalign{}
\endlastfoot
\textbf{Viewpoint} & \textbf{Stakeholder} & \textbf{Reason} \\
Direct & Construction Site Users (Builders, Owners, and Site
Supervisors) & Register construction sites, monitor progress, view
detected changes, review historical results, and generate reports. All
users have the same application permissions. \\
Indirect & Software Engineer & Develop, maintain, and update the
application and its backend services. \\
Indirect & AI Engineer & Develop, train, evaluate, and improve the
computer vision model and AI-based analysis. \\
Indirect & System Administrator & Maintain system availability, manage
technical configurations, and monitor system operation. \\
Domain & Google Earth Engine & Provide access to satellite imagery and
geospatial processing capabilities, subject to service limitations and
usage policies. \\
Domain & Sentinel-2 Satellite Imagery Provider & Supply satellite
imagery used for construction-site monitoring, subject to image
availability, spatial resolution, and data-usage conditions. \\
\end{longtable}

\renewcommand{\chapterlabel}{Chapter 3}
\chapter{3. Specific Requirements}\label{specific-requirements}

\section{3.1 External Interface
Requirements}\label{external-interface-requirements}

\subsection{3.1.1 User Interfaces}\label{user-interfaces}

The application interface should provide registration and login screens,
a dashboard, construction site registration with a map-based coordinate
confirmation step, site details and access-request screens, image
comparison controls, saved monitoring results, AI insights, and report
generation/download controls.

\subsection{3.1.2 Hardware Interfaces}\label{hardware-interfaces}

The user-side hardware requirements consist of a device capable of
running a modern web browser with a stable internet connection. The
server-side hardware requirements depend on the selected computer-vision
model and deployment environment and will be finalized during
implementation.

\subsection{3.1.3 Software Interfaces}\label{software-interfaces}

The system interfaces with Google Earth Engine for Sentinel-2 image
retrieval and geospatial processing; a database for user, site,
monitoring, and report metadata; a processing environment for image
preparation and change detection; and any configured AI component used
to generate observations. Exact database, map provider, AI service, and
storage implementation details should be recorded once finalized.

\subsection{3.1.4 Communications
Interfaces}\label{communications-interfaces}

The web application communicates with backend services over the internet
using standard web protocols, normally HTTPS for production traffic.
Backend services communicate with Google Earth Engine and the configured
database using their supported APIs or drivers. Authentication
credentials and sensitive data should be transmitted only through secure
channels.

\section{3.2 User Functional
Requirements}\label{user-functional-requirements}

\section{3. User Functional
Requirements}\label{user-functional-requirements-1}

\subsection{3.1. User Account
Management}\label{user-account-management}

\begin{enumerate}
\def\labelenumi{\arabic{enumi}.}
\item
  FR1: The system shall allow users to create an account by providing
  the required registration information.
\item
  FR2: The system shall allow registered users to log in using valid
  credentials.
\item
  FR3: The system shall allow authenticated users to log out of their
  accounts.
\end{enumerate}

\subsection{3.2. Data and Records
Management}\label{data-and-records-management}

\begin{enumerate}
\def\labelenumi{\arabic{enumi}.}
\item
  FR4: The system shall store and retrieve user information,
  construction site details, satellite imagery references, and
  monitoring results via database.
\item
  FR5: The system shall maintain records of monitoring operations and
  their associated results.
\item
  FR6: The system shall associate monitoring records with their
  corresponding construction sites.
\end{enumerate}

\subsection{3.3. Construction Change
Detection}\label{construction-change-detection}

\begin{enumerate}
\def\labelenumi{\arabic{enumi}.}
\item
  FR7: The system shall compare suitable satellite images captured at
  different times to identify changes within a registered construction
  site.
\item
  FR8: The system shall generate a change mask highlighting the areas
  where changes have been detected.
\item
  FR9: The system shall display the original satellite images alongside
  the corresponding change detection results.
\end{enumerate}

\subsection{3.4. AI-Based Construction
Analysis}\label{ai-based-construction-analysis}

\begin{enumerate}
\def\labelenumi{\arabic{enumi}.}
\item
  FR10: The system shall analyze detected changes to provide insights
  into construction site development over time.
\item
  FR11: The system shall summarize the detected changes in an
  understandable format for users.
\item
  FR12: The system shall present AI-generated observations based on the
  available image analysis results.
\end{enumerate}

\subsection{3.5. Dashboard and Results
Visualization}\label{dashboard-and-results-visualization}

\begin{itemize}
\item
  FR13: The system shall provide a dashboard displaying registered
  construction sites and their monitoring status.
\item
  FR14: The system shall display satellite images, detected changes, and
  calculated change measurements.
\item
  FR15: The system shall present monitoring results and detected changes
  through appropriate visual overlays.
\end{itemize}

\subsection{3.7. Construction Site Registration and
Management}\label{construction-site-registration-and-management}

\begin{itemize}
\item
  FR16: The system shall allow users to register a construction site by
  providing its name, geographical bounding-box coordinates, and
  construction start date.
\item
  FR17: The system shall validate the provided coordinates and display
  the selected location on a map for confirmation before completing
  registration.
\item
  FR18: The system shall allow users to view the details of their
  registered construction sites.
\end{itemize}

\subsection{\texorpdfstring{ 3.8. Construction Site Access and Member
Management}{ 3.8. Construction Site Access and Member Management}}\label{construction-site-access-and-member-management}

\begin{enumerate}
\def\labelenumi{\arabic{enumi}.}
\item
  FR19: The system shall allow users to request access to a construction
  site by providing its unique site identifier.
\item
  FR20: The system shall allow the creator of a construction site to
  approve or reject access requests from other users.
\end{enumerate}

\subsection{3.9. Construction Progress
Viewing}\label{construction-progress-viewing}

\begin{enumerate}
\def\labelenumi{\arabic{enumi}.}
\item
  FR21: The system shall allow users to view calculated change
  measurements for a registered construction site over time.
\item
  FR22: The system shall allow users to compare monitoring results from
  different dates to observe changes in the site.
\item
  FR23: The system shall allow users to view calculated change
  measurements for the latest satellite comparison of a registered
  construction site.
\end{enumerate}

\subsection{3.10. Monitoring Report
Generation}\label{monitoring-report-generation}

\begin{itemize}
\item
  FR24: The system shall allow users to generate reports for a selected
  construction site and time period.
\item
  FR25: The system shall include relevant satellite images, change
  detection results, and monitoring measurements in the generated
  reports.
\item
  FR26: The system shall allow users to download generated reports for
  record-keeping and review.
\end{itemize}

\subsection{3.11. Error Notifications and
Handling}\label{error-notifications-and-handling}

\begin{enumerate}
\def\labelenumi{\arabic{enumi}.}
\item
  FR27: The system shall notify users when suitable satellite imagery is
  unavailable for the selected site or date range.
\item
  FR28: The system shall inform users when a monitoring operation cannot
  proceed because of a processing or service error.
\item
  FR29: The system shall notify users when submitted construction site
  information is invalid or incomplete.
\item
  FR30: The system shall display an appropriate error message when a
  monitoring operation fails.
\item
  FR31: The system shall handle database connectivity failures and
  external service errors that prevent required application operations
  from completing.
\end{enumerate}

\textbf{Note:} All registered users will have the same access
permissions and functional capabilities within the application. The
functional requirements listed above describe the
system\textquotesingle s expected behavior from the
user\textquotesingle s perspective. These requirements will be further
refined into detailed System Functional Requirements (SFRs) to specify
the system\textquotesingle s operations and expected responses.

\section{3.3 System Functional
Requirements}\label{system-functional-requirements}

The following detailed system-level requirements are retained from the
supplied report. Each requirement table is preserved with its function
description, inputs and origin, outputs and destination, dependencies,
actions, preconditions, postconditions, and side effects where those
fields were present.

\section{1. User Account Management}\label{user-account-management-1}

This category covers account registration, login, and logout. We'll
begin with three system-level requirements.

\subsection{Requirement \# 1: User
Registration}\label{requirement-1-user-registration}

\begin{longtable}[]{@{}
  >{\raggedright\arraybackslash}p{(\columnwidth - 2\tabcolsep) * \real{0.1694}}
  >{\raggedright\arraybackslash}p{(\columnwidth - 2\tabcolsep) * \real{0.8306}}@{}}
\toprule\noalign{}
\endhead
\bottomrule\noalign{}
\endlastfoot
\textbf{Field} & \textbf{Description} \\
Function Description & The system shall allow a new user to create an
account by providing the required registration information. \\
Input \& Origin & User registration information, including name, email
address, and password, entered through the registration interface. \\
Outputs and Destination & A confirmation message indicating successful
registration, displayed to the user. The newly created account
information is stored in the database. If registration fails, an
appropriate error message is displayed. \\
Requires & Registration interface, user authentication component, and
database connectivity. \\
Actions to be taken & 1. Receive the registration information from the
user. 2. Validate that all required fields are provided and correctly
formatted. 3. Check whether the email address is already associated with
an existing account. 4. If the information is valid and the email is not
registered, securely process the password and create the account. 5.
Display the appropriate success or error message. \\
Pre-condition & The user is not logged in and has access to the
registration interface. \\
Post-condition & Upon successful registration, a new user account exists
in the database and can be used to access the system. \\
Side Effects & A new user record is created in the database. No
construction site or monitoring record is created at this stage. \\
\end{longtable}

\subsection{Requirement \# 2: User
Login}\label{requirement-2-user-login}

\begin{longtable}[]{@{}
  >{\raggedright\arraybackslash}p{(\columnwidth - 2\tabcolsep) * \real{0.1662}}
  >{\raggedright\arraybackslash}p{(\columnwidth - 2\tabcolsep) * \real{0.8338}}@{}}
\toprule\noalign{}
\endhead
\bottomrule\noalign{}
\endlastfoot
\textbf{Field} & \textbf{Description} \\
Function Description & The system shall allow registered users to log in
using their valid email address and password. \\
Input \& Origin & Email address and password entered by the user through
the login interface. \\
Outputs and Destination & Upon successful authentication, the user is
redirected to the main dashboard. If authentication fails, an
appropriate error message is displayed. \\
Requires & Login interface, user authentication component, and database
connectivity. \\
Actions to be taken & 1. Receive the email address and password. 2.
Validate that the required fields have been provided. 3. Retrieve the
corresponding account record using the email address. 4. Verify the
submitted password against the securely stored password information. 5.
If the credentials are valid, establish an authenticated session or
equivalent authentication state. 6. Redirect the user to the dashboard.
Otherwise, display an error message. \\
Pre-condition & The user has an existing account and is not currently
logged in. \\
Post-condition & Upon successful login, the user is authenticated and
can access the system\textquotesingle s available features. \\
Side Effects & An authenticated session or equivalent authentication
state is established. No construction site data is modified by the login
operation. \\
\end{longtable}

\subsection{Requirement \# 3: User
Logout}\label{requirement-3-user-logout}

\begin{longtable}[]{@{}
  >{\raggedright\arraybackslash}p{(\columnwidth - 2\tabcolsep) * \real{0.1711}}
  >{\raggedright\arraybackslash}p{(\columnwidth - 2\tabcolsep) * \real{0.8289}}@{}}
\toprule\noalign{}
\endhead
\bottomrule\noalign{}
\endlastfoot
\textbf{Field} & \textbf{Description} \\
Function Description & The system shall allow an authenticated user to
log out of their account. \\
Input \& Origin & A logout request initiated by the authenticated user
through the application interface. \\
Outputs and Destination & A logout confirmation or redirection to the
login page is provided to the user. \\
Requires & An active authenticated session or equivalent authentication
state and the logout functionality. \\
Actions to be taken & 1. Receive the logout request. 2. Invalidate the
current authenticated session or authentication state. 3. Clear relevant
authentication information maintained by the application. 4. Redirect
the user to the login page or another publicly accessible page. 5.
Ensure that protected pages and operations cannot be accessed using the
invalidated authentication state. \\
Pre-condition & The user is logged in to the system. \\
Post-condition & The user\textquotesingle s current authenticated
session is terminated, and access to protected features requires
authentication again. \\
Side Effects & The user may need to log in again to access protected
features. Existing construction site records and monitoring results
remain unchanged. \\
\end{longtable}

\section{2. Data and Records
Management}\label{data-and-records-management-1}

This category defines how the system stores, retrieves, and maintains
records related to users, registered construction sites, satellite
imagery, monitoring operations, and analysis results.

\subsection{Requirement \# 4: Store and Retrieve System
Records}\label{requirement-4-store-and-retrieve-system-records}

\begin{longtable}[]{@{}
  >{\raggedright\arraybackslash}p{(\columnwidth - 2\tabcolsep) * \real{0.1713}}
  >{\raggedright\arraybackslash}p{(\columnwidth - 2\tabcolsep) * \real{0.8287}}@{}}
\toprule\noalign{}
\endhead
\bottomrule\noalign{}
\endlastfoot
\textbf{Field} & \textbf{Description} \\
Function Description & The system shall store and retrieve the
information required to operate the construction monitoring system. \\
Input \& Origin & User account information from the registration
process; construction site details from the site registration interface;
satellite image metadata and references from GEE; monitoring results
from the computer vision and AI processing pipeline. \\
Outputs and Destination & Requested records are retrieved from the
database and provided to the relevant application component for display
or processing. New records are stored in the database. \\
Requires & Database, database access component, and the relevant
application modules. \\
Actions to be taken & 1. Receive data from the relevant application
component. 2. Validate the required data fields and formats. 3. Store
new records in the appropriate database tables. 4. Retrieve requested
records when needed by the application. 5. Return the retrieved records
to the requesting component. 6. Report an appropriate error if a
database operation fails. \\
Pre-condition & The database is configured and accessible, and the
requesting application component has supplied the required information
for the operation. \\
Post-condition & Valid information is stored successfully or the
requested records are returned to the requesting component. \\
Side Effects & Database records may be created, updated, or read. A
failed operation may leave the requested data unavailable until the
problem is resolved. \\
\end{longtable}

\subsection{Requirement \# 5: Maintain Monitoring Operation
Records}\label{requirement-5-maintain-monitoring-operation-records}

\begin{longtable}[]{@{}
  >{\raggedright\arraybackslash}p{(\columnwidth - 2\tabcolsep) * \real{0.1694}}
  >{\raggedright\arraybackslash}p{(\columnwidth - 2\tabcolsep) * \real{0.8306}}@{}}
\toprule\noalign{}
\endhead
\bottomrule\noalign{}
\endlastfoot
\textbf{Field} & \textbf{Description} \\
Function Description & The system shall maintain a record of each
monitoring operation and its associated processing status and
results. \\
Input \& Origin & Monitoring operation details from the monitoring
process management component, including the associated site, processing
dates, operation status, and available result references. \\
Outputs and Destination & A monitoring record containing the operation
details and current status is stored in the database and made available
to the relevant application components. \\
Requires & Database, monitoring process management component, and
computer vision processing pipeline. \\
Actions to be taken & 1. Receive the details of a monitoring operation.
2. Associate the operation with the relevant registered construction
site. 3. Create a monitoring record when the operation begins. 4. Update
the record as processing progresses. 5. Store the final status and
references to the generated results when processing finishes. 6. Record
a failure status if the operation cannot be completed. \\
Pre-condition & A construction site has been registered, and a
monitoring operation has been initiated or scheduled. \\
Post-condition & A record exists for the monitoring operation,
reflecting its latest known status and, when available, its results. \\
Side Effects & A new monitoring record may be created or an existing
record updated. The operation record itself does not guarantee that
suitable satellite imagery or successful analysis results are
available. \\
\end{longtable}

\subsection{Requirement \# 6: Associate Records with Their
Corresponding Construction
Sites}\label{requirement-6-associate-records-with-their-corresponding-construction-sites}

\begin{longtable}[]{@{}
  >{\raggedright\arraybackslash}p{(\columnwidth - 2\tabcolsep) * \real{0.1694}}
  >{\raggedright\arraybackslash}p{(\columnwidth - 2\tabcolsep) * \real{0.8306}}@{}}
\toprule\noalign{}
\endhead
\bottomrule\noalign{}
\endlastfoot
\textbf{Field} & \textbf{Description} \\
Function Description & The system shall associate satellite imagery
metadata, monitoring operations, and analysis results with their
corresponding registered construction sites. \\
Input \& Origin & Construction site identifier, satellite image
references and metadata from the imagery retrieval component, and
monitoring result references from the processing pipeline. \\
Outputs and Destination & Associated records are stored in the database
and can be retrieved by the relevant application components using the
construction site identifier. \\
Requires & Database, construction site records, imagery retrieval
component, and monitoring process management component. \\
Actions to be taken & 1. Receive the construction site identifier with
the relevant imagery, monitoring, or result information. 2. Verify that
the referenced construction site exists. 3. Associate the information
with that site using the appropriate database relationship. 4. Store the
association and related information. 5. Reject the association and
report an error if the site identifier is invalid. \\
Pre-condition & The construction site has been registered, and the
relevant imagery, monitoring, or result information is available for
association. \\
Post-condition & The relevant records are linked to the correct
construction site and can be retrieved for site monitoring only. \\
Side Effects & Database relationships may be created or updated. No
satellite image is downloaded or processed solely because a record is
associated with a site. \\
\end{longtable}

\section{3. Construction Change
Detection}\label{construction-change-detection-1}

This category defines how the system retrieves suitable Sentinel-2
satellite imagery through Google Earth Engine (GEE), prepares images for
comparison, detects changes between different dates, and presents the
detected areas.

\subsection{Requirement \# 7: Retrieve Suitable Satellite
Images}\label{requirement-7-retrieve-suitable-satellite-images}

\begin{longtable}[]{@{}
  >{\raggedright\arraybackslash}p{(\columnwidth - 2\tabcolsep) * \real{0.1645}}
  >{\raggedright\arraybackslash}p{(\columnwidth - 2\tabcolsep) * \real{0.8355}}@{}}
\toprule\noalign{}
\endhead
\bottomrule\noalign{}
\endlastfoot
\textbf{Field} & \textbf{Description} \\
Function Description & The system shall retrieve suitable satellite
imagery for a registered construction site and a specified monitoring
period using Google Earth Engine. \\
Input \& Origin & Construction site bounding-box coordinates and
monitoring date range from the database; satellite imagery availability
and metadata from GEE. \\
Outputs and Destination & References to suitable satellite images, their
acquisition dates, and relevant metadata are provided to the computer
vision processing pipeline and stored in the database. If no suitable
imagery is available, the system reports this condition to the
monitoring process management component. \\
Requires & Google Earth Engine access, registered construction site
information, database connectivity, and an imagery retrieval
component. \\
Actions to be taken & 1. Retrieve the selected site\textquotesingle s
coordinates and requested date range. 2. Query GEE for Sentinel-2
imagery covering the site and period. 3. Apply appropriate date, spatial
coverage, and image-quality filters, including cloud-related filtering.
4. Identify imagery suitable for comparison. 5. Return the image
references and relevant metadata to the processing pipeline. 6. Store
the references and metadata in the database. 7. Report when no suitable
imagery is found or retrieval fails. \\
Pre-condition & A construction site has been registered, a valid
monitoring period is available, and GEE access is configured. \\
Post-condition & Suitable image references and metadata are available
for processing, or the system records that suitable imagery could not be
retrieved. \\
Side Effects & GEE requests may consume available service quota or
processing resources. Imagery retrieval does not guarantee that every
requested date will have a usable image. \\
\end{longtable}

\subsection{Requirement \# 8: Preprocess Satellite Images for
Comparison}\label{requirement-8-preprocess-satellite-images-for-comparison}

\begin{longtable}[]{@{}
  >{\raggedright\arraybackslash}p{(\columnwidth - 2\tabcolsep) * \real{0.1648}}
  >{\raggedright\arraybackslash}p{(\columnwidth - 2\tabcolsep) * \real{0.8352}}@{}}
\toprule\noalign{}
\endhead
\bottomrule\noalign{}
\endlastfoot
\textbf{Field} & \textbf{Description} \\
Function Description & The system shall prepare retrieved satellite
images to make them suitable for change detection. \\
Input \& Origin & Satellite image pixel data, acquisition dates, and
image metadata retrieved through GEE or made available to the processing
environment. \\
Outputs and Destination & Preprocessed and comparable images are
provided to the computer vision change detection component. Relevant
preprocessing information may be stored with the monitoring record. \\
Requires & Satellite imagery retrieval component, image preprocessing
component, and access to the actual image pixel data. \\
Actions to be taken & 1. Retrieve or access the actual pixel data for
the selected images. 2. Apply appropriate quality and cloud masking
where required. 3. Ensure that images cover the registered construction
site. 4. Prepare images using consistent spatial coverage, resolution,
and data formats. 5. Check whether the resulting images are suitable for
comparison. 6. Pass suitable image pairs to the change detection
component or report why preprocessing could not be completed. \\
Pre-condition & At least two candidate satellite images are available
for the selected site and monitoring period. \\
Post-condition & Suitable preprocessed image pairs are available for
change detection, or the system identifies that the images cannot be
compared reliably. \\
Side Effects & Temporary image data and intermediate processing files
may be created. Preprocessing may exclude images or areas that do not
meet the required quality conditions. \\
\end{longtable}

\subsection{Requirement \# 9: Detect Changes Between Satellite
Images}\label{requirement-9-detect-changes-between-satellite-images}

\begin{longtable}[]{@{}
  >{\raggedright\arraybackslash}p{(\columnwidth - 2\tabcolsep) * \real{0.1662}}
  >{\raggedright\arraybackslash}p{(\columnwidth - 2\tabcolsep) * \real{0.8338}}@{}}
\toprule\noalign{}
\endhead
\bottomrule\noalign{}
\endlastfoot
\textbf{Field} & \textbf{Description} \\
Function Description & The system shall compare suitable satellite
images acquired at different times and identify areas where detectable
changes have occurred within a registered construction site. \\
Input \& Origin & A suitable pair of preprocessed satellite images from
the image preprocessing component and the corresponding site information
from the database. \\
Outputs and Destination & A change mask identifying detected changed and
unchanged areas, along with relevant processing metadata, is provided to
the results visualization and analysis components. \\
Requires & Preprocessed image pairs, a trained or configured computer
vision change detection model, and the processing environment. \\
Actions to be taken & 1. Receive the selected image pair and associated
metadata. 2. Provide the images to the change detection model. 3.
Compare the images to identify areas exhibiting detectable changes. 4.
Generate a change mask representing the detected changes. 5. Associate
the output with the corresponding site and image dates. 6. Pass the
change mask and processing information to the relevant components. 7.
Report an error if inference fails. \\
Pre-condition & Two suitable, preprocessed satellite images of the same
site from different dates are available, and the change detection model
is ready for use. \\
Post-condition & A change mask and associated metadata are generated and
made available to the results processing components, or a failure status
is recorded. \\
Side Effects & Processing consumes computing resources. Detection errors
may occur because of image quality, resolution, seasonal differences,
shadows, or model limitations. \\
\end{longtable}

\subsection{Requirement \# 10: Calculate and Store Detected Change
Measurements}\label{requirement-10-calculate-and-store-detected-change-measurements}

\begin{longtable}[]{@{}
  >{\raggedright\arraybackslash}p{(\columnwidth - 2\tabcolsep) * \real{0.1629}}
  >{\raggedright\arraybackslash}p{(\columnwidth - 2\tabcolsep) * \real{0.8371}}@{}}
\toprule\noalign{}
\endhead
\bottomrule\noalign{}
\endlastfoot
\textbf{Field} & \textbf{Description} \\
Function Description & The system shall calculate the estimated area
represented by detected changes and associate the measurements with the
corresponding monitoring operation. \\
Input \& Origin & Generated change mask, image resolution or spatial
metadata, site boundary information, and monitoring dates from the
change detection component and database. \\
Outputs and Destination & Estimated changed-area measurements,
measurement units, and associated dates are stored in the database and
provided to the dashboard and AI-based analysis components. \\
Requires & Change mask, geospatial metadata, area calculation component,
and database connectivity. \\
Actions to be taken & 1. Receive the generated change mask and relevant
spatial metadata. 2. Identify the pixels classified as changed within
the site boundary. 3. Calculate the estimated area represented by the
detected pixels using the appropriate spatial resolution or geospatial
calculation method. 4. Associate the measurement with the relevant image
dates and monitoring operation. 5. Store the measurement and its unit in
the database. 6. Report an error if the required spatial information is
unavailable or the calculation fails. \\
Pre-condition & A valid change mask and sufficient spatial metadata are
available for the registered construction site. \\
Post-condition & The estimated changed area is calculated and stored
with its corresponding monitoring record, or the system records that the
measurement could not be calculated. \\
Side Effects & New measurement records may be created. The calculated
area represents detected image changes and does not necessarily equal
the actual construction area or percentage of construction completed. \\
\end{longtable}

\subsection{Requirement \# 11: Display Change Detection
Results}\label{requirement-11-display-change-detection-results}

\begin{longtable}[]{@{}
  >{\raggedright\arraybackslash}p{(\columnwidth - 2\tabcolsep) * \real{0.1678}}
  >{\raggedright\arraybackslash}p{(\columnwidth - 2\tabcolsep) * \real{0.8322}}@{}}
\toprule\noalign{}
\endhead
\bottomrule\noalign{}
\endlastfoot
\textbf{Field} & \textbf{Description} \\
Function Description & The system shall display the satellite images
used for comparison and the corresponding detected changes to the
user. \\
Input \& Origin & Original or processed satellite images, acquisition
dates, generated change masks, and calculated change measurements from
the database and processing pipeline. \\
Outputs and Destination & A visual comparison of the selected images,
highlighted changed areas, acquisition dates, and available change
measurements is displayed on the application\textquotesingle s results
interface. \\
Requires & Results visualization component, image and result references,
and access to the corresponding image data and generated outputs. \\
Actions to be taken & 1. Receive the selected monitoring result. 2.
Retrieve the corresponding satellite images or their stored
representations. 3. Retrieve the associated change mask and
measurements. 4. Display the images with their acquisition dates. 5.
Overlay or otherwise visualize the detected changed areas. 6. Display
the available measurements and report an error if the required results
cannot be retrieved. \\
Pre-condition & At least one completed change detection operation has
produced results available for display. \\
Post-condition & The user can view the relevant images, detected
changes, and available measurements for the selected monitoring
operation. \\
Side Effects & Image and result data may be loaded into the
user\textquotesingle s browser or application session. Displaying
results does not modify the stored analysis records. \\
\end{longtable}

\section{4. AI-Based Construction
Analysis}\label{ai-based-construction-analysis-1}

This category defines how the system interprets the outputs of computer
vision analysis and presents understandable observations about changes
detected at a construction site.

\subsection{Requirement \# 12: Analyze Detected Construction
Changes}\label{requirement-12-analyze-detected-construction-changes}

\begin{longtable}[]{@{}
  >{\raggedright\arraybackslash}p{(\columnwidth - 2\tabcolsep) * \real{0.1662}}
  >{\raggedright\arraybackslash}p{(\columnwidth - 2\tabcolsep) * \real{0.8338}}@{}}
\toprule\noalign{}
\endhead
\bottomrule\noalign{}
\endlastfoot
\textbf{Field} & \textbf{Description} \\
Function Description & The system shall analyze detected change
measurements and associated monitoring information to identify patterns
in construction site development over time. \\
Input \& Origin & Change masks, estimated changed areas, satellite image
acquisition dates, and historical monitoring results from the computer
vision pipeline and database. \\
Outputs and Destination & Analysis findings, identified change patterns,
and relevant observations are provided to the dashboard and AI-generated
insights component. \\
Requires & Completed change detection results, stored monitoring
records, and an analysis component. \\
Actions to be taken & 1. Retrieve the available change measurements and
corresponding monitoring dates. 2. Identify observable patterns, such as
increases or decreases in detected changed areas. 3. Generate analysis
findings based on the available evidence. 4. Make the findings available
to the dashboard and insights component. \\
Pre-condition & At least one completed monitoring operation has produced
valid change measurements. \\
Post-condition & Analysis findings are generated from the available
monitoring results and associated with the relevant construction site
and monitoring period. \\
Side Effects & Additional analysis records may be stored. Incomplete or
inconsistent monitoring data may limit the reliability of identified
patterns. \\
\end{longtable}

\subsection{Requirement \# 13: Generate AI-Based Construction
Insights}\label{requirement-13-generate-ai-based-construction-insights}

\begin{longtable}[]{@{}
  >{\raggedright\arraybackslash}p{(\columnwidth - 2\tabcolsep) * \real{0.1629}}
  >{\raggedright\arraybackslash}p{(\columnwidth - 2\tabcolsep) * \real{0.8371}}@{}}
\toprule\noalign{}
\endhead
\bottomrule\noalign{}
\endlastfoot
\textbf{Field} & \textbf{Description} \\
Function Description & The system shall generate understandable AI-based
observations from available construction monitoring results to help
users interpret detected changes. \\
Input \& Origin & Change detection measurements, image dates, and
relevant analysis findings supplied by the computer vision pipeline and
database. \\
Outputs and Destination & A textual summary of observed changes,
possible development patterns, and relevant limitations is provided to
the dashboard. \\
Requires & Available monitoring results, analysis findings, and an
AI-based insight generation component. \\
Actions to be taken & 1. Retrieve the relevant monitoring results and
analysis findings. 2. Prepare the available measurements, and dates as
input to the AI component. 3. Generate a concise explanation of the
observed changes. 4. Ensure that the explanation is based on the
supplied evidence and does not present unsupported conclusions as facts.
5. Associate the generated insights with the corresponding site and
monitoring period. 6. Provide the insights to the dashboard. 7. Report
an appropriate error if insight generation fails. \\
Pre-condition & Relevant monitoring results or analysis findings are
available, and the AI component is configured and accessible. \\
Post-condition & AI-generated observations are available for display, or
the system records that insight generation was unsuccessful. \\
Side Effects & AI processing may consume computing resources or external
service quota, depending on the implementation. Generated observations
may contain inaccuracies and must not be treated as definitive evidence
of construction completion or structural safety. \\
\end{longtable}

\subsection{}\label{section}

\subsection{}\label{section-1}

\subsection{}\label{section-2}
\newpage
\subsection{Requirement \# 14: Display AI Analysis and Its Supporting
Measurements}\label{requirement-14-display-ai-analysis-and-its-supporting-measurements}

\begin{longtable}[]{@{}
  >{\raggedright\arraybackslash}p{(\columnwidth - 2\tabcolsep) * \real{0.1694}}
  >{\raggedright\arraybackslash}p{(\columnwidth - 2\tabcolsep) * \real{0.8306}}@{}}
\toprule\noalign{}
\endhead
\bottomrule\noalign{}
\endlastfoot
\textbf{Field} & \textbf{Description} \\
Function Description & The system shall display AI-generated
construction insights alongside the corresponding monitoring
measurements and dates. \\
Input \& Origin & Generated AI observations, estimated changed-area
measurements, monitoring dates, and for a selected monitoring check and
allow users to generate and download reports. from the analysis
component and database. \\
Outputs and Destination & An integrated presentation of AI insights and
supporting monitoring data is displayed on the construction
site\textquotesingle s results dashboard. \\
Requires & AI insight generation component, stored monitoring results,
database connectivity, and dashboard visualization component. \\
Actions to be taken & 1. Retrieve the AI-generated observations for the
selected site or monitoring period. 2. Retrieve the measurements and
dates associated with those observations. 3. Display the observations
alongside their supporting information. 4. Identify the relevant
monitoring period for each observation. 5. Indicate when AI insights are
unavailable or could not be generated. \\
Pre-condition & The user has selected a registered construction site or
monitoring record, and the relevant results are available. \\
Post-condition & The user can review the generated observations together
with the measurements and monitoring dates used to support them. \\
Side Effects & Displaying the analysis does not modify the underlying
satellite imagery, change masks, or monitoring measurements. \\
\end{longtable}

\section{5. Dashboard and Results
Visualization}\label{dashboard-and-results-visualization-1}

This category defines how the system presents registered construction
sites, monitoring operations, satellite imagery, detected changes,
measurements, and AI-generated insights through a user-friendly
dashboard.

\subsection{Requirement \# 15: Display the Main
Dashboard}\label{requirement-15-display-the-main-dashboard}

\begin{longtable}[]{@{}
  >{\raggedright\arraybackslash}p{(\columnwidth - 2\tabcolsep) * \real{0.1664}}
  >{\raggedright\arraybackslash}p{(\columnwidth - 2\tabcolsep) * \real{0.8336}}@{}}
\toprule\noalign{}
\endhead
\bottomrule\noalign{}
\endlastfoot
\textbf{Field} & \textbf{Description} \\
Function Description & The system shall provide an interactive dashboard
displaying the user\textquotesingle s registered construction sites and
their monitoring status. \\
Input \& Origin & Registered construction site details and monitoring
status records retrieved from the database. \\
Outputs and Destination & A dashboard containing the registered site
list, site details, and monitoring statuses is displayed to the user. \\
Requires & User authentication, database connectivity, construction site
management component, and dashboard interface. \\
Actions to be taken & 1. Verify that the user is authenticated. 2.
Retrieve the construction sites associated with the
user\textquotesingle s account. 3. Retrieve the latest available
monitoring status for each site. 4. Display the site names and relevant
summary information. 5. Allow the user to select a site to view its
details and monitoring results. 6. Display an appropriate message if no
sites have been registered or if the information cannot be retrieved. \\
Pre-condition & The user is authenticated and the dashboard interface is
accessible. \\
Post-condition & The dashboard displays the available construction site
information and monitoring statuses, or an appropriate empty-state or
error message. \\
Side Effects & Dashboard data may be retrieved from the database. No
site details or monitoring results are modified by displaying the
dashboard. \\
\end{longtable}

\subsection{Requirement \# 16: Display Construction Site Monitoring
Results}\label{requirement-16-display-construction-site-monitoring-results}

\begin{longtable}[]{@{}
  >{\raggedright\arraybackslash}p{(\columnwidth - 2\tabcolsep) * \real{0.1681}}
  >{\raggedright\arraybackslash}p{(\columnwidth - 2\tabcolsep) * \real{0.8319}}@{}}
\toprule\noalign{}
\endhead
\bottomrule\noalign{}
\endlastfoot
\textbf{Field} & \textbf{Description} \\
Function Description & The system shall display the available monitoring
results for a selected construction site. \\
Input \& Origin & Selected construction site identifier, satellite image
references, acquisition dates, change masks, estimated changed areas,
and AI-generated insights from the database and processing
components. \\
Outputs and Destination & A site-specific results view displaying
available satellite images, detected changes, measurements, monitoring
dates, and AI-generated observations. \\
Requires & Registered construction site information, completed
monitoring results, database connectivity, and results visualization
component. \\
Actions to be taken & 1. Receive the user\textquotesingle s selection of
a construction site. 2. Retrieve the associated monitoring records. 3.
Retrieve the available image and change detection output references. 4.
Display the relevant images, acquisition dates, and detected changes. 5.
Display calculated measurements and available AI-generated insights. 6.
Indicate when results are unavailable or could not be retrieved. \\
Pre-condition & The user is authenticated and has selected a registered
construction site. \\
Post-condition & The available monitoring results for the selected site
are displayed, or the system informs the user that no results are
available. \\
Side Effects & Relevant images and result data may be loaded into the
application session. Stored monitoring records remain unchanged. \\
\end{longtable}

\subsection{Requirement \# 17: Display Construction Site Location on
a Map}\label{requirement-17-display-construction-site-location-on-a-map}

\begin{longtable}[]{@{}
  >{\raggedright\arraybackslash}p{(\columnwidth - 2\tabcolsep) * \real{0.1711}}
  >{\raggedright\arraybackslash}p{(\columnwidth - 2\tabcolsep) * \real{0.8289}}@{}}
\toprule\noalign{}
\endhead
\bottomrule\noalign{}
\endlastfoot
\textbf{Field} & \textbf{Description} \\
Function Description & The system shall display the geographical
location and registered bounding box of a selected construction site on
a map. \\
Input \& Origin & Construction site bounding-box coordinates and site
name retrieved from the database. \\
Outputs and Destination & A map displaying the selected construction
site\textquotesingle s location and geographical boundary is presented
to the user. \\
Requires & Registered site coordinates, map visualization component, and
access to the map data or map service used by the application. \\
Actions to be taken & 1. Receive the selected construction site
identifier. 2. Retrieve its stored coordinates. 3. Convert the
coordinates into the format required by the map component. 4. Display
the location and bounding box on the map. 5. Adjust the map view to show
the selected site. 6. Display an appropriate message if the coordinates
are unavailable or cannot be displayed. \\
Pre-condition & The user is authenticated and has selected a registered
construction site with valid geographical coordinates. \\
Post-condition & The site\textquotesingle s registered location and
boundary are displayed on the map, or the system reports that the
location cannot be displayed. \\
Side Effects & The map component may request map data from an external
service. Displaying the location does not modify the registered
coordinates. \\
\end{longtable}

\section{7. Construction Site Registration and
Management}\label{construction-site-registration-and-management-1}

This category defines how users register construction sites, validate
geographical coordinates, confirm the selected location on a map, and
view their registered site information.

\subsection{Requirement \# 18: Register a Construction
Site}\label{requirement-23-register-a-construction-site}

\begin{longtable}[]{@{}
  >{\raggedright\arraybackslash}p{(\columnwidth - 2\tabcolsep) * \real{0.1629}}
  >{\raggedright\arraybackslash}p{(\columnwidth - 2\tabcolsep) * \real{0.8371}}@{}}
\toprule\noalign{}
\endhead
\bottomrule\noalign{}
\endlastfoot
\textbf{Field} & \textbf{Description} \\
Function Description & The system shall allow an authenticated user to
register a construction site by providing its name, geographical
bounding-box coordinates, and construction start date. \\
Input \& Origin & Site name, four bounding-box coordinates (minimum and
maximum latitude and longitude), and construction start date entered
through the site registration interface. \\
Outputs and Destination & Upon successful registration, a confirmation
message is displayed, and the construction site information is stored in
the database. If registration fails, an appropriate error message is
displayed. \\
Requires & User authentication, site registration interface, coordinate
validation component, database connectivity, and map visualization
component. \\
Actions to be taken & \begin{minipage}[t]{\linewidth}\raggedright
1. Receive the site name, bounding-box coordinates, and construction
start date.\\
2. Validate that all required fields are provided.\\
3. Validate the coordinate values and ensure that the bounding box is
geographically valid. 4. Verify that the construction start date is
valid and not in the future.\\
5. Display the selected location on a map for user confirmation.\\
6. After confirmation, store the site information in the database and
associate it with the registering user.\\
7. Display a success message or an appropriate validation or storage
error.\strut
\end{minipage} \\
Pre-condition & The user is authenticated and has access to the
construction site registration interface. \\
Post-condition & A valid, confirmed construction site is stored in the
database and associated with the registering user, or registration
remains incomplete if validation, confirmation, or storage fails. \\
Side Effects & A new construction site record may be created. The site
becomes available for subsequent monitoring operations after successful
registration. \\
\end{longtable}

\subsection{Requirement \# 19: Validate and Confirm Construction Site
Coordinates}\label{requirement-24-validate-and-confirm-construction-site-coordinates}

\begin{longtable}[]{@{}
  >{\raggedright\arraybackslash}p{(\columnwidth - 2\tabcolsep) * \real{0.1613}}
  >{\raggedright\arraybackslash}p{(\columnwidth - 2\tabcolsep) * \real{0.8387}}@{}}
\toprule\noalign{}
\endhead
\bottomrule\noalign{}
\endlastfoot
\textbf{Field} & \textbf{Description} \\
Function Description & The system shall validate the submitted
geographical coordinates and display the corresponding bounding box on a
map before completing site registration. \\
Input \& Origin & Minimum and maximum latitude and longitude entered by
the user through the site registration interface. \\
Outputs and Destination & Coordinate validation results and a map
displaying the selected geographical boundary are presented to the user.
Invalid coordinates result in an appropriate error message. \\
Requires & Site registration interface, coordinate validation component,
and map visualization component. \\
Actions to be taken & \begin{minipage}[t]{\linewidth}\raggedright
1. Receive the four bounding-box coordinates.\\
2. Check that all coordinate values are numeric and within valid
latitude and longitude ranges.\\
3. Verify that the minimum latitude is less than the maximum latitude
and the minimum longitude is less than the maximum longitude.\\
4. If the coordinates are valid, display the bounding box on the map.\\
5. Ask the user to confirm that the displayed location is correct.\\
6. Allow the user to correct the coordinates if the location is
incorrect.\\
7. Prevent registration from proceeding with invalid or unconfirmed
coordinates.\strut
\end{minipage} \\
Pre-condition & The user has opened the site registration interface and
entered the geographical coordinates. \\
Post-condition & Valid coordinates are displayed for confirmation, or
the user is informed that the coordinates are invalid. Site registration
can proceed only after successful validation and confirmation. \\
Side Effects & The map component may request map data from an external
service. No construction site record is finalized solely by validating
or displaying the coordinates. \\
\end{longtable}

\subsection{}\label{section-3}

\subsection{Requirement \# 20: View Registered Construction Site
Details}\label{requirement-25-view-registered-construction-site-details}

\begin{longtable}[]{@{}
  >{\raggedright\arraybackslash}p{(\columnwidth - 2\tabcolsep) * \real{0.1694}}
  >{\raggedright\arraybackslash}p{(\columnwidth - 2\tabcolsep) * \real{0.8306}}@{}}
\toprule\noalign{}
\endhead
\bottomrule\noalign{}
\endlastfoot
\textbf{Field} & \textbf{Description} \\
Function Description & The system shall allow authenticated users to
view the details of their registered construction sites. \\
Input \& Origin & A selected construction site identifier or a request
to view registered sites submitted through the dashboard. \\
Outputs and Destination & The site\textquotesingle s name, geographical
coordinates, construction start date, and other available site
information are displayed to the user. \\
Requires & User authentication, construction site records, database
connectivity, and site details interface. \\
Actions to be taken & \begin{minipage}[t]{\linewidth}\raggedright
1. Receive the request to view registered sites or a specific
site\textquotesingle s details.\\
2. Retrieve the relevant site records associated with the
user\textquotesingle s account.\\
3. If a specific site is selected, retrieve its stored details.\\
4. Display the available information in the site details interface.\\
5. Display an appropriate message if no sites are registered or the
requested record cannot be retrieved.\strut
\end{minipage} \\
Pre-condition & The user is authenticated and has access to the site
management interface. \\
Post-condition & The requested registered site information is displayed,
or the system informs the user that the information is unavailable. \\
Side Effects & Site records may be retrieved from the database. Viewing
site details does not modify the stored information. \\
\end{longtable}

\section{8. Construction Site Registration and
Management}\label{construction-site-registration-and-management-2}

\subsection{Requirement \# 21: Request Access to a Construction Site}\label{sfr-21-request-access}

\begin{longtable}[]{@{}
  >{\raggedright\arraybackslash}p{(\columnwidth - 2\tabcolsep) * \real{0.1680}}
  >{\raggedright\arraybackslash}p{(\columnwidth - 2\tabcolsep) * \real{0.8320}}@{}}
\toprule\noalign{}
\begin{minipage}[b]{\linewidth}\raggedright
Field
\end{minipage} & \begin{minipage}[b]{\linewidth}\raggedright
Description
\end{minipage} \\
\midrule\noalign{}
\endhead
\bottomrule\noalign{}
\endlastfoot
Function Description & The system shall allow users to request access to
a registered construction site by providing its unique site
identifier. \\
Input \& Origin & Unique construction site identifier entered by the
user and validated against construction site records stored in the
database. \\
Outputs and Destination & A confirmation message indicating successful
submission, or an appropriate message explaining why the request could
not be submitted, shall be displayed to the user. The request details
shall be stored in the database for review by the construction site
creator. \\
Requires & User authentication and an existing registered construction
site. \\
Actions to be Taken & Verify that the site exists and check whether the
user already has access or has an existing pending request. If no
duplicate request exists, create an access request associated with the
user and the construction site. \\
Pre-condition & The user must be logged in, and the construction site
must be registered in the system. \\
Post-condition & A valid access request is stored with a pending status
for review by the construction site creator. \\
Side Effects & A new access request record is added to the database. \\
\end{longtable}

\subsection{Requirement \# 22: Approve or Reject Construction Site Access Requests}\label{sfr-22-approve-reject-access}

\begin{longtable}[]{@{}
  >{\raggedright\arraybackslash}p{(\columnwidth - 2\tabcolsep) * \real{0.1651}}
  >{\raggedright\arraybackslash}p{(\columnwidth - 2\tabcolsep) * \real{0.8349}}@{}}
\toprule\noalign{}
\begin{minipage}[b]{\linewidth}\raggedright
Field
\end{minipage} & \begin{minipage}[b]{\linewidth}\raggedright
Description
\end{minipage} \\
\midrule\noalign{}
\endhead
\bottomrule\noalign{}
\endlastfoot
Function Description & The system shall allow the creator of a
construction site to review access requests and approve or reject
requests submitted by other users. \\
Input \& Origin & Construction site identifier, access request details,
and the approval or rejection decision selected by the construction site
creator through the user interface. \\
Outputs and Destination & The updated request status and the result of
the decision shall be displayed to the construction site creator. The
access request status shall be updated in the database, and access
permissions shall be updated when the request is approved. \\
Requires & User authentication, construction site records, and pending
access requests. \\
Actions to be Taken & Verify that the user making the decision is the
creator of the relevant construction site and that the request is
pending. If approved, grant the requesting user access and update the
request status to approved. If rejected, update the request status to
rejected without granting access. \\
Pre-condition & The user must be logged in as the creator of the
relevant construction site, and the access request must have a pending
status. \\
Post-condition & The access request status is updated to approved or
rejected. Access is granted only when the request is approved. \\
Side Effects & The access request record and, when applicable, the
site\textquotesingle s membership and access permissions are updated in
the database. \\
\end{longtable}

\section{8. Construction Progress
Monitoring}\label{construction-progress-monitoring}

This category defines how the system presents monitoring results over
time, compares observations from different dates, and helps users assess
changes in construction sites using available satellite imagery and
calculated measurements.

\subsection{Requirement \# 23: View Construction Change Measurements
Over
Time}\label{requirement-26-view-construction-change-measurements-over-time}

\begin{longtable}[]{@{}
  >{\raggedright\arraybackslash}p{(\columnwidth - 2\tabcolsep) * \real{0.1678}}
  >{\raggedright\arraybackslash}p{(\columnwidth - 2\tabcolsep) * \real{0.8322}}@{}}
\toprule\noalign{}
\endhead
\bottomrule\noalign{}
\endlastfoot
\textbf{Field} & \textbf{Description} \\
Function Description & The system shall allow users to view calculated
changed-area measurements for a registered construction site viewing the
latest calculated change measurement. \\
Input \& Origin & Selected construction site identifier and available
monitoring records containing acquisition dates and estimated
changed-area measurements, retrieved from the database. \\
Outputs and Destination & A chronological list or visualization of
available changed-area measurements, including their corresponding
monitoring dates and measurement units, is displayed to the user. \\
Requires & Registered construction site information, historical
monitoring records, database connectivity, and dashboard visualization
component. \\
Actions to be taken & \begin{minipage}[t]{\linewidth}\raggedright
1. Receive the selected construction site identifier.\\
2. Retrieve the site\textquotesingle s available monitoring measurements
and associated dates.\\
3. Arrange the measurements in chronological order.\\
4. Display each measurement with its corresponding date and unit.\\
5. Indicate when measurements are unavailable for some monitoring
periods.\\
6. Display an appropriate message if no valid measurements are
available.\strut
\end{minipage} \\
Pre-condition & The user is authenticated and has selected a registered
construction site. \\
Post-condition & Available change measurements are displayed
chronologically, or the system informs the user that no valid
measurements are available. \\
Side Effects & Historical monitoring records may be retrieved for
display. Stored measurements remain unchanged. \\
\end{longtable}

\section{9. Monitoring Report
Generation}\label{monitoring-report-generation-1}

This category defines how users generate and download reports containing
construction site monitoring information, satellite imagery, detected
changes, measurements, and relevant AI-generated insights.

\subsection{Requirement \# 24: Generate a Construction Monitoring
Report}\label{requirement-29-generate-a-construction-monitoring-report}

\begin{longtable}[]{@{}
  >{\raggedright\arraybackslash}p{(\columnwidth - 2\tabcolsep) * \real{0.1613}}
  >{\raggedright\arraybackslash}p{(\columnwidth - 2\tabcolsep) * \real{0.8387}}@{}}
\toprule\noalign{}
\endhead
\bottomrule\noalign{}
\endlastfoot
\textbf{Field} & \textbf{Description} \\
Function Description & The system shall allow users to generate a
monitoring report for a selected construction site and monitoring
period. \\
Input \& Origin & Selected construction site, requested date range, and
report generation request submitted by the user through the
dashboard. \\
Outputs and Destination & A monitoring report containing the available
information for the selected site and period is generated and made
available to the user. \\
Requires & Registered construction site information, monitoring records,
report generation component, database connectivity, and access to the
relevant monitoring results. \\
Actions to be taken & \begin{minipage}[t]{\linewidth}\raggedright
1. Receive the selected construction site and requested date range.\\
2. Validate the selected site and date range.\\
3. Retrieve the relevant monitoring records within the requested
period.\\
4. Collect the available change detection results, measurements, and
AI-generated insights. 5. Organize the information into a structured
report.\\
6. Generate the report in pdf.\\
7. Make the generated report available for download.\\
8. Display an appropriate message if the report cannot be generated or
the requested period contains no usable results.\strut
\end{minipage} \\
Pre-condition & The user is authenticated, a construction site is
selected, and the requested reporting period is valid. \\
Post-condition & A report is generated for the selected site and period
when sufficient information is available, or the system informs the user
that report generation was unsuccessful. \\
Side Effects & Report files may be created in temporary or persistent
storage. Report generation may consume processing resources. \\
\end{longtable}

\subsection{Requirement \# 25: Include Monitoring Information in the
Report}\label{requirement-30-include-monitoring-information-in-the-report}

\begin{longtable}[]{@{}
  >{\raggedright\arraybackslash}p{(\columnwidth - 2\tabcolsep) * \real{0.1613}}
  >{\raggedright\arraybackslash}p{(\columnwidth - 2\tabcolsep) * \real{0.8387}}@{}}
\toprule\noalign{}
\endhead
\bottomrule\noalign{}
\endlastfoot
\textbf{Field} & \textbf{Description} \\
Function Description & The system shall include the relevant
construction site and monitoring information available for the selected
reporting period in the generated report. \\
Input \& Origin & Construction site details, monitoring dates, estimated
changed-area measurements, and AI-generated insights retrieved from the
database and processing pipeline. \\
Outputs and Destination & A structured report containing the applicable
site information, monitoring observations, and supporting visual or
numerical results is provided to the report generation component. \\
Requires & Site records, completed monitoring results, image and output
references, report generation component, and database connectivity. \\
Actions to be taken & \begin{minipage}[t]{\linewidth}\raggedright
1. Retrieve the registered construction site details.\\
2. Retrieve monitoring records corresponding to the selected period.\\
3. Include the site name, geographical information, and reporting
period.\\
4. Include acquisition dates.\\
5. Include the corresponding change detection visualizations and
estimated measurements. 6. Include available AI-generated observations
and relevant limitations.\\
7. Ensure that each result is associated with its corresponding
monitoring date.\\
8. Pass the assembled information to the report generation
component.\strut
\end{minipage} \\
Pre-condition & A valid report request exists, and relevant site
information and monitoring records are available. \\
Post-condition & The report contains the available information relevant
to the selected site and reporting period, with results associated with
their corresponding dates. \\
Side Effects & Selected images and analysis results may be retrieved or
temporarily assembled for report generation. The original records remain
unchanged. \\
\end{longtable}

\subsection{Requirement \# 26: Download a Generated Monitoring
Report}\label{requirement-31-download-a-generated-monitoring-report}

\begin{longtable}[]{@{}
  >{\raggedright\arraybackslash}p{(\columnwidth - 2\tabcolsep) * \real{0.1726}}
  >{\raggedright\arraybackslash}p{(\columnwidth - 2\tabcolsep) * \real{0.8274}}@{}}
\toprule\noalign{}
\endhead
\bottomrule\noalign{}
\endlastfoot
\textbf{Field} & \textbf{Description} \\
Function Description & The system shall allow users to download a
successfully generated construction monitoring report. \\
Input \& Origin & A download request for a generated report, initiated
by the user through the report interface. \\
Outputs and Destination & The generated report file is delivered to the
user\textquotesingle s browser for downloading to their device. \\
Requires & A successfully generated report, report file storage or
retrieval mechanism, user authentication, and report download
functionality. \\
Actions to be taken & \begin{minipage}[t]{\linewidth}\raggedright
1. Receive the user\textquotesingle s download request.\\
2. Identify the requested report.\\
3. Create a Report\\
4. Deliver the file to the user\textquotesingle s browser in pdf
format.\\
6. Display an appropriate error message if the report is unavailable or
the download fails.\strut
\end{minipage} \\
Pre-condition & The user is authenticated, and a successfully generated
report is available for download. \\
Post-condition & The report file is made available to the user for
downloading, or the system informs the user that the download could not
be completed. \\
Side Effects & Network bandwidth may be consumed during the download.
Downloading a report does not modify the stored monitoring results. \\
\end{longtable}

\section{10. Error Notifications and
Handling}\label{error-notifications-and-handling-1}

This category defines how the system identifies errors, informs users
about problems, and handles situations that prevent site registration,
satellite imagery retrieval, monitoring, analysis, or report generation
from completing successfully.

\subsection{Requirement \# 27: Notify Users About Invalid or
Incomplete Site
Information}\label{requirement-32-notify-users-about-invalid-or-incomplete-site-information}

\begin{longtable}[]{@{}
  >{\raggedright\arraybackslash}p{(\columnwidth - 2\tabcolsep) * \real{0.1662}}
  >{\raggedright\arraybackslash}p{(\columnwidth - 2\tabcolsep) * \real{0.8338}}@{}}
\toprule\noalign{}
\endhead
\bottomrule\noalign{}
\endlastfoot
\textbf{Field} & \textbf{Description} \\
Function Description & The system shall notify users when submitted
construction site information is incomplete or invalid. \\
Input \& Origin & Site name, geographical bounding-box coordinates,
construction start date, and other required site information submitted
through the site registration interface. \\
Outputs and Destination & An appropriate validation error message is
displayed to the user, identifying the field or information that
requires correction. \\
Requires & Site registration interface, input validation component, and
coordinate validation component. \\
Actions to be taken & 1. Receive the submitted construction site
information. 2. Check that all required fields are provided. 3. Validate
the field formats and coordinate ranges. 4. Verify that the bounding-box
coordinates define a valid geographical area. 5. Validate the
construction start date. 6. Identify any invalid or missing information.
7. Display an appropriate error message and prevent registration until
the required corrections are made. \\
Pre-condition & The user has opened the site registration interface and
submitted or attempted to submit construction site information. \\
Post-condition & Invalid or incomplete information is identified and
communicated to the user. The site is not registered until the required
validation checks are passed. \\
Side Effects & Validation messages may be displayed, but no valid
construction site record is created from the rejected submission. \\
\end{longtable}

\subsection{Requirement \# 28: Notify Users When Suitable Satellite
Imagery Is
Unavailable}\label{requirement-33-notify-users-when-suitable-satellite-imagery-is-unavailable}

\begin{longtable}[]{@{}
  >{\raggedright\arraybackslash}p{(\columnwidth - 2\tabcolsep) * \real{0.1645}}
  >{\raggedright\arraybackslash}p{(\columnwidth - 2\tabcolsep) * \real{0.8355}}@{}}
\toprule\noalign{}
\endhead
\bottomrule\noalign{}
\endlastfoot
\textbf{Field} & \textbf{Description} \\
Function Description & The system shall notify users when suitable
satellite imagery cannot be found for the selected construction site and
monitoring period. \\
Input \& Origin & Construction site coordinates, requested monitoring
date range, and imagery availability or quality information returned by
GEE. \\
Outputs and Destination & A notification explaining that suitable
satellite imagery is unavailable is displayed to the user and associated
with the relevant monitoring operation. \\
Requires & GEE access, imagery retrieval component, image-quality
filtering component, and monitoring process management component. \\
Actions to be taken & 1. Receive the imagery retrieval results. 2.
Determine whether suitable images are available for the requested site
and period. 3. Identify whether the issue results from unavailable
imagery, image quality, cloud cover, or insufficient suitable image
pairs, where this information is available. 4. Record the appropriate
monitoring status. 5. Display an understandable notification to the
user. 6. Allow the user to select another monitoring period or retry
when appropriate. \\
Pre-condition & A monitoring operation has requested satellite imagery
for a registered construction site. \\
Post-condition & The user is informed when suitable imagery is
unavailable, and the monitoring operation is recorded with an
appropriate status. \\
Side Effects & The monitoring operation may remain incomplete.
Previously completed monitoring results remain available. \\
\end{longtable}

\subsection{Requirement \# 29: Handle Monitoring and Processing
Errors}\label{requirement-34-handle-monitoring-and-processing-errors}

\begin{longtable}[]{@{}
  >{\raggedright\arraybackslash}p{(\columnwidth - 2\tabcolsep) * \real{0.1662}}
  >{\raggedright\arraybackslash}p{(\columnwidth - 2\tabcolsep) * \real{0.8338}}@{}}
\toprule\noalign{}
\endhead
\bottomrule\noalign{}
\endlastfoot
\textbf{Field} & \textbf{Description} \\
Function Description & The system shall detect and handle errors
occurring during satellite imagery retrieval, image preprocessing,
change detection, measurement calculation, or AI analysis. \\
Input \& Origin & Error codes, exception information, service responses,
and processing status information generated by the relevant system
components. \\
Outputs and Destination & An appropriate error notification is displayed
to the user, and relevant failure information is recorded in the
monitoring operation record. \\
Requires & Error handling mechanism, monitoring process management
component, processing pipeline, database connectivity, and logging
facility. \\
Actions to be taken & 1. Receive an error from the affected component.
2. Identify the associated monitoring operation and processing stage. 3.
Record the failure status and relevant diagnostic information. 4.
Prevent the operation from being incorrectly marked as successfully
completed. 5. Preserve previously completed results. 6. Display a
user-friendly notification without exposing sensitive technical details.
7. Permit a retry when appropriate and safe. \\
Pre-condition & A monitoring or analysis operation is in progress, and a
system component reports an error. \\
Post-condition & The error is handled, relevant failure information is
recorded, and the user is informed of the unsuccessful operation or
affected processing stage. \\
Side Effects & Partial outputs or temporary processing files may remain
and require cleanup. Retrying an operation may consume additional
computing resources or GEE service quota. \\
\end{longtable}
\newpage
\subsection{Requirement \# 30: Handle Report Generation and Download
Errors}\label{requirement-35-handle-report-generation-and-download-errors}

\begin{longtable}[]{@{}
  >{\raggedright\arraybackslash}p{(\columnwidth - 2\tabcolsep) * \real{0.1662}}
  >{\raggedright\arraybackslash}p{(\columnwidth - 2\tabcolsep) * \real{0.8338}}@{}}
\toprule\noalign{}
\endhead
\bottomrule\noalign{}
\endlastfoot
\textbf{Field} & \textbf{Description} \\
Function Description & The system shall notify users when a monitoring
report cannot be generated or downloaded successfully. \\
Input \& Origin & Report generation requests, report data availability
information, file storage responses, and download status information
from the report generation and download components. \\
Outputs and Destination & An appropriate error message is displayed to
the user, and relevant diagnostic information is recorded where
applicable. \\
Requires & Report generation component, report storage or retrieval
mechanism, database connectivity, and error handling mechanism. \\
Actions to be taken & 1. Receive an error from the report generation or
download component. 2. Identify the affected report operation. 3.
Determine whether the failure occurred during data retrieval, report
creation, file storage, or file delivery. 4. Record the
operation\textquotesingle s failure status and relevant diagnostic
information. 5. Display an appropriate message explaining that the
report is unavailable or the operation failed. 6. Allow the user to
retry when appropriate and safe. \\
Pre-condition & A user has requested a monitoring report to be generated
or downloaded. \\
Post-condition & The user is informed of the failure, and the system
records the relevant operation status where applicable. If the failure
is resolved and the operation is retried successfully, the report
becomes available. \\
Side Effects & Incomplete report files or temporary files may require
cleanup. A retry may consume additional processing resources. \\
\end{longtable}

\subsection{Requirement \# 31: Handle Database and External Service
Errors}\label{requirement-36-handle-database-and-external-service-errors}

\begin{longtable}[]{@{}
  >{\raggedright\arraybackslash}p{(\columnwidth - 2\tabcolsep) * \real{0.1645}}
  >{\raggedright\arraybackslash}p{(\columnwidth - 2\tabcolsep) * \real{0.8355}}@{}}
\toprule\noalign{}
\endhead
\bottomrule\noalign{}
\endlastfoot
\textbf{Field} & \textbf{Description} \\
Function Description & The system shall handle database connectivity
failures and external service errors that prevent required application
operations from completing. \\
Input \& Origin & Database error responses, GEE service errors,
authentication or access failures, network errors, and other relevant
service status information. \\
Outputs and Destination & An appropriate error status is recorded where
possible, and a user-friendly notification is displayed when the failure
affects a requested operation. \\
Requires & Database access component, external service integration
components, error handling mechanism, and application logging
facility. \\
Actions to be taken & 1. Detect a database or external service failure.
2. Identify the affected operation and service. 3. Record the error
details when logging or database access remains available. 4. Prevent
incomplete operations from being reported as successful. 5. Display an
appropriate message to the user. 6. Retry transient failures only when
appropriate and safe. 7. Resume or restart the affected operation when
the service becomes available and the operation can safely continue. \\
Pre-condition & An application operation requires database access or an
external service, and the required resource becomes unavailable or
returns an error. \\
Post-condition & The failure is handled and communicated appropriately.
The operation either completes successfully after recovery or remains
marked as failed or incomplete. \\
Side Effects & Some operations may be delayed or temporarily
unavailable. Error logging may be incomplete if the database or logging
service itself is unavailable. Repeated retries may consume additional
resources. \\
\end{longtable}

\subsection{Requirement \# 32: Handle Monitoring Operation
Failures}\label{requirement-22-handle-monitoring-operation-failures}

\begin{longtable}[]{@{}
  >{\raggedright\arraybackslash}p{(\columnwidth - 2\tabcolsep) * \real{0.1645}}
  >{\raggedright\arraybackslash}p{(\columnwidth - 2\tabcolsep) * \real{0.8355}}@{}}
\toprule\noalign{}
\endhead
\bottomrule\noalign{}
\endlastfoot
\textbf{Field} & \textbf{Description} \\
Function Description & The system shall identify failures that prevent a
monitoring operation from completing successfully. \\
Input \& Origin & Error information from GEE requests, image
preprocessing, computer vision inference, measurement calculation,
database operations, or other monitoring components. \\
Outputs and Destination & The system shall display the
operation\textquotesingle s failure status and details on the user
interface. \\
Requires & Monitoring process management component, error handling
mechanism, database connectivity, and access to operation status
records. \\
Actions to be taken & 1. Receive an error reported by a monitoring
component. 2. Identify the affected monitoring operation.4. Prevent the
failed operation from being incorrectly marked as successfully
completed. 5. Display a user-appropriate error message. 6. Preserve
previously completed monitoring results. 7. Allow a new attempt when the
cause has been addressed and the operation can safely be retried. \\
Pre-condition & A monitoring operation has been initiated, and a
component reports an error or failure. \\
Post-condition & The failure is communicated appropriately, while
previously completed results remain available. \\
Side Effects & Partial processing outputs or temporary files may remain
and may require cleanup. A retry may consume additional resources or
external service quota. \\
\end{longtable}
\newpage
\section{3.4 Behaviour Requirements}\label{behaviour-requirements}

This section contains the interaction and behavioural models associated
with the requirements above.

\subsection{3.4.1 Use Case Diagram}\label{use-case-diagram}

\begin{center}
    \IfFileExists{Use_Case_Diagram.png}
    {\includegraphics[width=0.8\textwidth]{Use_Case_Diagram.png}}
    {\textit{Use Case Diagram image not found.}}
\end{center}

\subsection{3.4.2 Data Flow Diagram --- Level 0}
\label{data-flow-diagram-level-0}

\begin{center}
    \IfFileExists{DFD_Level_0.png}
    {\includegraphics[width=0.8\textwidth]{DFD_Level_0.png}}
    {\textit{DFD Level 0 image not found.}}
\end{center}

\subsection{3.4.3 Data Flow Diagram --- Level 1}
\label{data-flow-diagram-level-1}

\begin{center}
    \IfFileExists{DFD_Level_1.png}
    {\includegraphics[width=0.8\textwidth]{DFD_Level_1.png}}
    {\textit{DFD Level 1 image not found.}}
\end{center}

\subsection{3.4.4 Data Flow Diagram --- Level 2: Compare Images}
\label{data-flow-diagram-level-2-compare-images}

\begin{center}
    \IfFileExists{DFD_Level_2_Compare_Images.png}
    {\includegraphics[width=0.8\textwidth]{DFD_Level_2_Compare_Images.png}}
    {\textit{DFD Level 2: Compare Images diagram not found.}}
\end{center}

\subsection{3.4.5 Sequence Diagram}
\label{sequence-diagram}

\begin{center}
    \IfFileExists{Sequence.png}
    {\includegraphics[width=0.9\textwidth]{Sequence.png}}
    {\textit{Sequence Diagram image not found.}}
\end{center}

\renewcommand{\chapterlabel}{Chapter 4}
\chapter{4. Other Non-Functional Requirements}
\label{other-non-functional-requirements}

\section{Non-Functional Requirements (NFRs)}

\section{1. Performance and Response
Time}\label{performance-and-response-time}

\begin{longtable}[]{@{}
  >{\raggedright\arraybackslash}p{(\columnwidth - 4\tabcolsep) * \real{0.1089}}
  >{\raggedright\arraybackslash}p{(\columnwidth - 4\tabcolsep) * \real{0.2142}}
  >{\raggedright\arraybackslash}p{(\columnwidth - 4\tabcolsep) * \real{0.6769}}@{}}
\toprule\noalign{}
\endhead
\bottomrule\noalign{}
\endlastfoot
\textbf{Req ID} & \textbf{Requirement Name} & \textbf{Non-Functional
Requirement Specification} \\
NFR-PERF-01 & Dashboard Load Time & The system shall load and display
registered construction sites within 3 seconds. \\
NFR-PERF-02 & Image Processing Time & The system shall complete
satellite image preprocessing and change detection within 120 seconds
per bounding box for the defined test area, number of image pairs, and
hardware setup. If processing takes longer, it shall continue in the
background and show the current status to the user. \\
NFR-PERF-03 & Report Generation Time & The system shall generate a PDF
monitoring report within 10 seconds, provided that the required
monitoring data and images are already available. \\
NFR-PERF-04 & User Interface Responsiveness & The system shall provide
visible feedback within 1 seconds when a user performs an action, such
as selecting a date, entering coordinates, or changing a visualization
layer. \\
\end{longtable}

\section{2. Accuracy and Quality}\label{accuracy-and-quality}

\begin{longtable}[]{@{}
  >{\raggedright\arraybackslash}p{(\columnwidth - 4\tabcolsep) * \real{0.1057}}
  >{\raggedright\arraybackslash}p{(\columnwidth - 4\tabcolsep) * \real{0.2158}}
  >{\raggedright\arraybackslash}p{(\columnwidth - 4\tabcolsep) * \real{0.6784}}@{}}
\toprule\noalign{}
\endhead
\bottomrule\noalign{}
\endlastfoot
\textbf{Req ID} & \textbf{Requirement Name} & \textbf{Non-Functional
Requirement Specification} \\
NFR-ACC-01 & Change Detection Model Quality & The computer vision model
shall achieve an Intersection over Union (IoU) score of at least 0.50 on
a separate test dataset containing manually labelled satellite images
representative of the project\textquotesingle s target imagery. \\
NFR-ACC-02 & Spatial Resolution & The system shall process selected
Sentinel-2 bands at their native 10-metre spatial resolution, where
applicable. It shall not claim that the detected changes provide more
spatial detail than the source imagery supports. \\
NFR-ACC-03 & Input Validation & The system shall reject all invalid
inputs included in the defined validation test cases, including missing
required fields, coordinates outside valid ranges, incorrect bounding
boxes, and invalid construction start dates. \\
\end{longtable}

\section{3. Availability and
Reliability}\label{availability-and-reliability}

\begin{longtable}[]{@{}
  >{\raggedright\arraybackslash}p{(\columnwidth - 4\tabcolsep) * \real{0.1026}}
  >{\raggedright\arraybackslash}p{(\columnwidth - 4\tabcolsep) * \real{0.2126}}
  >{\raggedright\arraybackslash}p{(\columnwidth - 4\tabcolsep) * \real{0.6848}}@{}}
\toprule\noalign{}
\endhead
\bottomrule\noalign{}
\endlastfoot
\textbf{Req ID} & \textbf{Requirement Name} & \textbf{Non-Functional
Requirement Specification} \\
NFR-REL-01 & System Availability & The web application and backend API
shall remain available for at least 95\% of each calendar month,
excluding scheduled maintenance. \\
NFR-REL-02 & External Service Error Handling & If Google Earth Engine
(GEE) or another external service fails, the system shall display an
appropriate error message within 2 seconds of detecting the failure and
record the status of the affected operation whenever possible. \\
\end{longtable}

\phantomsection\label{_rh6k7hor8fco}{}

\section{4. Capacity and Scalability}\label{capacity-and-scalability}

\begin{longtable}[]{@{}
  >{\raggedright\arraybackslash}p{(\columnwidth - 4\tabcolsep) * \real{0.1009}}
  >{\raggedright\arraybackslash}p{(\columnwidth - 4\tabcolsep) * \real{0.1967}}
  >{\raggedright\arraybackslash}p{(\columnwidth - 4\tabcolsep) * \real{0.7024}}@{}}
\toprule\noalign{}
\endhead
\bottomrule\noalign{}
\endlastfoot
\textbf{Req ID} & \textbf{Requirement Name} & \textbf{Non-Functional
Requirement Specification} \\
NFR-SCA-01 & Construction Site Capacity & The system shall support at
least 1,000 registered construction sites without requiring changes to
the database structure. \\
NFR-SCA-02 & Processed Image Caching & The system shall temporarily
store eligible processed image outputs for up to 10 days to reduce
repeated processing and unnecessary requests to external services.
Cached results shall be updated when their source data or processing
settings change. \\
\end{longtable}
\vspace{1cm}
\section{5. Security and Compliance}\label{security-and-compliance}

\begin{longtable}[]{@{}
  >{\raggedright\arraybackslash}p{(\columnwidth - 4\tabcolsep) * \real{0.1074}}
  >{\raggedright\arraybackslash}p{(\columnwidth - 4\tabcolsep) * \real{0.2189}}
  >{\raggedright\arraybackslash}p{(\columnwidth - 4\tabcolsep) * \real{0.6737}}@{}}
\toprule\noalign{}
\endhead
\bottomrule\noalign{}
\endlastfoot
\textbf{Req ID} & \textbf{Requirement Name} & \textbf{Non-Functional
Requirement Specification} \\
NFR-SEC-01 & Session Inactivity Timeout & The system shall automatically
end an authenticated session after 1 hour of inactivity. The user shall
log in again to access protected features. \\
NFR-SEC-02 & Password Length & The system shall require passwords to
contain at least 8 characters and shall store them using Django's
password-hashing method instead of storing them as plain text. \\
\end{longtable}

\renewcommand{\chapterlabel}{Chapter 5}
\chapter{5. Other Requirements}\label{other-requirements}

No additional requirements beyond those specified in Chapters 2--4 and
Appendix A have been identified at this stage.

\renewcommand{\chapterlabel}{}
\chapter{Appendix A --- Data
Dictionary}\label{appendix-a-data-dictionary}

\textbf{Data Model and Database Specification:}

Data Dictionary

\section{1.1 users Table}\label{users-table}

This table stores the account information of registered users.

\begin{quote}
· \textbf{user\_id} (UUID, Primary Key): A unique identifier assigned to
each user.

· \textbf{email} (VARCHAR, Unique, Not Null): The user\textquotesingle s
email address, used for account identification.

· \textbf{password\_hash} (VARCHAR, Not Null): The securely hashed
password, if password management is handled by the application.

· \textbf{full\_name} (VARCHAR): The full name of the user.
\end{quote}

\section{1.2 construction\_sites
Table}\label{construction_sites-table}

This table stores the details of construction sites registered for
monitoring.

\begin{quote}
· \textbf{construction\_id} (UUID, Primary Key): A unique identifier
assigned to each construction site.

· \textbf{site\_name} (VARCHAR, Not Null): The name of the construction
site.

· \textbf{created\_by} (UUID, Foreign Key → users.user\_id): Identifies
the user who registered the site.

· \textbf{bounding\_boundaries} (JSON) The coordinates pointing to the
boundaries of the whole construction site.\\
\textbf{location} (VARCHAR): A general address or description of the
site\textquotesingle s location.

· \textbf{city} (VARCHAR): The city in which the site is located.

· \textbf{country} (VARCHAR): The country in which the site is located.

· \textbf{area} (DECIMAL): The approximate area of the construction site
in square metres.

· \textbf{start\_date} (DATE, Not Null): The construction start date
provided by the user.

· \textbf{is\_completed} (BOOLEAN, Default: False): Indicates whether
the user has marked the construction project as completed.

· \textbf{end\_date} (DATE, Nullable): The actual completion date, if
the project has been completed.

· \textbf{is\_public} (BOOLEAN, Default: False): Indicates whether the
site is publicly viewable.

· \textbf{created\_at} (TIMESTAMP): The date and time when the site
record was created.

· \textbf{current\_progress} (DECIMAL): Latest calculated overall
progress (0.00\% to 100.00\%).
\end{quote}

\section{1.3 user\_site Table}\label{user_site-table}

This table links users to construction sites they are permitted to
access. It supports a many-to-many relationship because a user may
access multiple sites, and a site may be associated with multiple users.

\begin{quote}
· \textbf{User\_site\_id} (UUID, Primary Key): A unique identifier for
the access record.

· \textbf{user\_id} (UUID, Foreign Key → users.user\_id): Identifies the
user receiving access.

· \textbf{construction\_id} (UUID, Foreign Key →
construction\_sites.construction\_id): Identifies the construction site.

· \textbf{granted\_at} (TIMESTAMP): The date and time when access was
granted.
\end{quote}

\section{1.4 site\_access\_request
Table}\label{site_access_request-table}

This table stores requests from users who want to access private
construction sites.

\begin{quote}
· \textbf{request\_id} (UUID, Primary Key): A unique identifier for each
request.

· \textbf{request\_from\_id} (UUID, Foreign Key → users.user\_id):
Identifies the user requesting access.

· \textbf{request\_to\_id} (UUID, Foreign Key → users.user\_id):
Identifies the user receiving the request, usually the
site\textquotesingle s owner.

· \textbf{construction\_id} (UUID, Foreign Key →
construction\_sites.construction\_id): Identifies the site for which
access is requested.

· \textbf{status} (VARCHAR, Default:
\textquotesingle PENDING\textquotesingle): Indicates whether the request
is pending, approved, or rejected.

· \textbf{created\_at} (TIMESTAMP): The date and time when the request
was submitted.
\end{quote}

\section{1.5 progress\_logs Table}\label{progress_logs-table}

This table stores the latest run output for the active site check.
performed on a construction site. Each record represents a comparison
between satellite images taken at different times.

\begin{enumerate}
\def\labelenumi{\arabic{enumi}.}
\item
  \textbf{progress\_id} (UUID, Primary Key): A unique identifier for
  each monitoring operation.
\item
  \textbf{construction\_id} (UUID, Foreign Key →
  construction\_sites.construction\_id): Identifies the construction
  site being monitored.
\item
  \textbf{Baseline\_image\_file\_name} (VARCHAR, Nullable): The filename
  of the original baseline satellite image stored in Supabase Storage
  (e.g., construction\_id\_\_baseline\_date.tif).
\item
  \textbf{latest\_image\_filename (VARCHAR}, Nullable\textbf{): The
  filename of the latest satellite image stored in Supabase Storage
  using the standard format:
  \textless construction\_id\textgreater\_\_\textless date\_of\_satellite\_image\textgreater.tif
  (e.g., c123\_456\_\_2026-10-19.tif).}
\item
  \textbf{change\_mask\_filename (VARCHAR}, Nullable\textbf{): The
  filename of the OpenCV-generated visual overlay mask stored in
  Supabase Storage (e.g., c123\_456\_\_2026-10-19\_mask.tif).}
\item
  \textbf{area\_changed\_sqm} (DECIMAL, Nullable): The estimated area of
  detected change, in square metres.
\item
  \textbf{change\_percentage} (DECIMAL, Nullable): The percentage of the
  monitored site area classified as changed, if calculated by the
  system.
\item
  \textbf{status} (VARCHAR): The current status of the monitoring
  operation, such as PENDING, PROCESSING, COMPLETED, or FAILED.
\item
  \textbf{Satellite\_image\_date} (DATE): Date the satellite image was
  captured.
\item
  \textbf{processed\_at} (TIMESTAMP, Nullable): The date and time when
  the monitoring operation finished processing.
\end{enumerate}

\begin{quote}
·
\end{quote}

\subsection{2. ENTITY RELATIONSHIPS}\label{entity-relationships}

The database tables are connected through primary and foreign keys to
maintain data consistency and establish relationships between users,
construction sites, access requests, and monitoring records.

\textbf{1. Users and Construction Sites (One-to-Many, 1)}

One user can register multiple construction sites, while each
construction site is registered by one user. This relationship is
established through the created\_by field in the construction\_sites
table, which references users.user\_id.

\textbf{2. Users and Construction Sites through User\_Site
(Many-to-Many, M)}

A user can have access to multiple construction sites, and a
construction site can be accessed by multiple users. This many-to-many
relationship is managed through the user\_site table, which connects
users and construction\_sites using user\_id and construction\_id.

\textbf{3. Users and Site Access Requests (One-to-Many, 1)}

A user can submit multiple access requests to view private construction
sites. This relationship is established through the request\_from\_id
field in the site\_access\_request table, which references
users.user\_id.

A user can also receive multiple access requests from other users. This
relationship is established through the request\_to\_id field, which
also references users.user\_id.

\textbf{4. Construction Sites and Site Access Requests (One-to-Many, 1)}

A construction site can receive multiple access requests from different
users. Each access request is associated with one construction site
through the construction\_id field in the site\_access\_request table,
which references construction\_sites.construction\_id.

\textbf{5. Construction Sites and Progress Logs (One-to-Many, 1)}

A construction site can have multiple progress logs, with each log
representing a separate monitoring operation. This relationship is
established through the construction\_id field in the progress\_logs
table, which references construction\_sites.construction\_id.

These relationships allow the system to maintain site information,
manage user access, process access requests, and store the latest
monitoring result for each construction site.
\newpage
\chapter{Appendix B --- Requirements Traceability
Matrix}\label{appendix-b-requirements-traceability-matrix}

\begin{longtable}[]{@{}
  >{\raggedright\arraybackslash}p{(\columnwidth - 8\tabcolsep) * \real{0.0641}}
  >{\raggedright\arraybackslash}p{(\columnwidth - 8\tabcolsep) * \real{0.2878}}
  >{\raggedright\arraybackslash}p{(\columnwidth - 8\tabcolsep) * \real{0.1034}}
  >{\raggedright\arraybackslash}p{(\columnwidth - 8\tabcolsep) * \real{0.3580}}
  >{\raggedright\arraybackslash}p{(\columnwidth - 8\tabcolsep) * \real{0.1867}}@{}}
\toprule\noalign{}
\begin{minipage}[b]{\linewidth}\raggedright
\textbf{UFR ID}
\end{minipage} & \begin{minipage}[b]{\linewidth}\raggedright
\textbf{User Requirement / Summary}
\end{minipage} & \begin{minipage}[b]{\linewidth}\raggedright
\textbf{Related SFR ID(s)}
\end{minipage} & \begin{minipage}[b]{\linewidth}\raggedright
\textbf{Acceptance Criteria / Verification}
\end{minipage} & \begin{minipage}[b]{\linewidth}\raggedright
\textbf{Use Case / Model}
\end{minipage} \\
\midrule\noalign{}
\endhead
\bottomrule\noalign{}
\endlastfoot
FR1 & Create a user account. & SFR1 & Valid registration details create
an account and display a confirmation message. & Sign Up \\
FR2 & Log in using valid credentials. & SFR2 & Valid credentials allow
login; invalid credentials are rejected. & Login \\
FR3 & Log out of the account. & SFR3 & Logging out ends the session and
prevents access to protected pages until login. & Logout \\
FR4 & Store and retrieve user, site, imagery, and monitoring records. &
SFR4 & Saved records can be retrieved with their stored information
intact. & Database \\
FR5 & Maintain monitoring operation records and results. & SFR5 & A
monitoring operation\textquotesingle s status and available results are
saved and retrievable. & Compare Images; View Monitoring Results \\
FR6 & Associate monitoring records with their corresponding sites. &
SFR6 & Each monitoring record is linked to the correct construction
site. & View Monitoring Results \\
FR7 & Compare satellite images captured at different times. & SFR7; SFR8; SFR9 &
Given suitable images of the same site from different dates, the system
compares them and identifies detectable changes. & Compare Images \\
FR8 & Generate a change mask highlighting detected areas. & SFR10 & A
successful change detection operation produces a mask highlighting
detected changes. & Compare Images \\
FR9 & Display original satellite images alongside change detection
results. & SFR11 & The comparison images and corresponding change mask
are displayed with the correct dates. & View Monitoring Results \\
FR10 & Analyze detected changes to provide construction insights. &
SFR12 & Available change detection results are analyzed and produce
relevant findings. & AI Analysis \\
FR11 & Summarize detected changes understandably. & SFR13 & The system
displays a readable summary based on available monitoring results. & AI
Analysis \\
FR12 & Present AI-generated observations based on image analysis. &
SFR14 & AI observations are displayed with supporting monitoring data
and do not present unsupported conclusions as facts. & AI Analysis \\
FR13 & Display registered sites and monitoring status on a dashboard. &
SFR15 & The dashboard displays the user\textquotesingle s accessible
sites and their available monitoring statuses. & View Dashboard \\
FR14 & Display satellite images, detected changes, and measurements. &
SFR16 & Available images, detected changes, and calculated measurements
are displayed for the selected monitoring result. & View Monitoring
Results \\
FR15 & Present monitoring results through visual overlays. & SFR17 & The
system overlays the generated change mask on the corresponding satellite
imagery. & View Monitoring Results \\
FR16 & Register a site with its name, coordinates, and construction
start date. & SFR18 & Valid site details are saved, and the system
confirms successful registration. & Register Construction Site \\
FR17 & Validate coordinates and show the location for confirmation. &
SFR19 & Valid coordinates show the selected location; invalid
coordinates are rejected with a message. & Register Construction Site \\
FR18 & View registered construction site details. & SFR20 & Selecting a
site displays its saved details correctly. & View Construction Site \\
FR19 & Request access using a unique site identifier. & SFR21 & A user
without access can submit a request for an existing site; duplicate
pending requests are rejected. & Request to Join Construction Site \\
FR20 & Allow the site creator to approve or reject access requests. &
SFR22 & The creator can approve or reject a pending request; only
approval grants access. & Manage Site Access Requests \\
FR21 & View calculated change measurements over time. & SFR23 & The
system displays available measurements with their corresponding dates
and units. & View Monitoring Results \\
FR22 & Compare monitoring results from different dates. & SFR23 &
Selecting different monitoring dates displays the corresponding results
for comparison. & Compare Images; View Monitoring Results \\
FR23 & View measurements from the latest satellite comparison. & SFR23 &
The latest successfully completed comparison\textquotesingle s
measurements are displayed for the selected site. & View Monitoring
Results \\
FR24 & Generate a report for a selected site and time period. & SFR24 &
Selecting a site and a valid period with available results produces a
report. & Report Generation \\
FR25 & Include images, change detection results, and measurements in
reports. & SFR25 & The report includes relevant available imagery,
change detection results, measurements, and dates. & Report
Generation \\
FR26 & Download generated monitoring reports. & SFR26 & The user can
download and open a successfully generated report. & Report
Generation \\
FR27 & Notify users when suitable satellite imagery is unavailable. &
SFR27 & If no suitable imagery exists for the selected site or period,
the system displays a notification. & Compare Images \\
FR28 & Inform users when monitoring cannot proceed because of a
processing or service error. & SFR28 & When a required service or
processing component fails, the system displays an appropriate error
message. & Compare Images\\
FR29 & Notify users when site information is invalid or incomplete. &
SFR28 & Missing required fields or invalid site information prevent
registration and produce a validation message. & Register Construction
Site \\
FR30 & Display an error message during failure. & SFR29; SFR30; SFR31 SFR32 & A
failed monitoring, report generation, database or construction site operation is marked as unsuccessful, the user is
informed, and previously completed results remain available. & Compare
Images, Report Generation, Database, View Monitoring Results\\
\end{longtable}

\end{document}